% Unit 130 English translation of chapter9.tex, lines 1457--1625.
% Source: Methods of Algebra, Volume 2, commit
% 9a5803ff77dd3257484cb177f851a73770a59dd3 (CC BY 4.0).
% stable-unit-id: o014.aljabr2.chapter9.hopf-reconstruction
% Coverage: Section 9.6, reconstruction of coalgebras, bialgebras, and the
% antipode.
% Verified source corrections: the lower arrow on source line 1574 is
% omega({}^*f), as required by its source and target; source line 1578 starts
% from [phi tensor x], as in the preceding definition of S.

\section{Reconstruction of Hopf Algebras}\label{sec:Hopf-reconstruction}
This section continues the discussion of \S\ref{sec:Tannaka-coend}. Theorem
\ref{prop:Tannaka-aux1} constructs the coalgebra $\coEnd(\omega)$ from the
data $(\mathcal{A}, \omega)$ and identifies $\mathcal{A}$ with
$\catesubd{Comod}{\mathrm{f}}\coEnd(\omega)$. The most basic situation is
when $B = \Bbbk$ is a field. For a $\Bbbk$-coalgebra $L$, take
\begin{gather*}
	\mathcal{A} := \catesubd{Comod}{\mathrm{f}}L , \\
	\omega := U: \catesubd{Comod}{\mathrm{f}}L \xrightarrow{\text{forget}} \cate{Vect}_{\mathrm{f}}(\Bbbk).
\end{gather*}

Since $\Bbbk$ is a field, Proposition \ref{prop:Comod-abelian} with
$B = \Bbbk$ shows that $\cated{Comod}L$ is an abelian category, while
$\catesubd{Comod}{\mathrm{f}}L$, formed by the finite-dimensional comodules,
is clearly a locally finite abelian subcategory. It is now natural to ask:
does the corresponding $\coEnd(\omega)$ reconstruct the original datum $L$,
or does it instead produce a new coalgebra? As expected, the answer is
``reconstruction'', but a little argument is required.

The first step is to make clear how to compare $\coEnd(\omega)$ and $L$.
Since $L$ coacts from the right on every $\omega(X)$, the universal property
induces a coalgebra homomorphism
\begin{equation}\label{eqn:u-comparison}
	u: \coEnd(\omega) \to L.
\end{equation}
The original question is equivalent to asking whether $u$ is an isomorphism.
We first prove a simple result.

\begin{lemma}\label{prop:comodule-fd-union}
	Let $\Bbbk$ be a field and let $L$ be a $\Bbbk$-coalgebra, with the
	comultiplication homomorphism denoted by $\Delta$ and the counit
	homomorphism by $\epsilon$. Every right $L$-comodule $V$ is a union of
	finite-dimensional subcomodules.
\end{lemma}
\begin{proof}
	Suppose that the comodule structure is given by
	$\rho: V \to V \dotimes{\Bbbk} L$. For each
	$x \in V \dotimes{\Bbbk} L$, define the subspace
	\[ V(x) := \lrangle{(\identity \otimes \lambda)(x) : \lambda \in \Hom_{\Bbbk}(L, \Bbbk) } \; \subset V ; \]
	if $x$ is written concretely as
	$\sum_{i=1}^n v_i \otimes \ell_i$, with
	$\ell_1, \ldots, \ell_n$ linearly independent, then clearly
	$V(x) = \lrangle{v_1, \ldots, v_n}$, and hence
	$x \in V(x) \dotimes{\Bbbk} L$.
	
	Take $v \in V$. Setting $x = \rho(v)$ and $\lambda = \epsilon$ gives
	$v \in V(\rho(v))$. We next show that
	$\rho(V(\rho(v))) \subset V(\rho(v)) \dotimes{\Bbbk} L$. Write again
	$\rho(v) = \sum_i v_i \otimes \ell_i$; the equation
	$(\rho \otimes \identity)\rho(v) = (\identity \otimes \Delta)\rho(v)$
	in the definition of a comodule can be written as the following equation
	in $V \dotimes{\Bbbk} L \dotimes{\Bbbk} L$:
	\[ \sum_i \rho(v_i) \otimes \ell_i = \sum_i v_i \otimes \Delta(\ell_i). \]
	
	Extend $\ell_1, \ldots, \ell_n$ to a basis and contract with the dual
	basis. This gives $\rho(v_i) \in V(\rho(v)) \dotimes{\Bbbk} L$ for
	$i = 1, \ldots, n$.
	
	Thus every finite-dimensional subspace $\lrangle{v_1, \ldots, v_m}$ of
	$V$ is contained in the subcomodule
	$\sum_{i=1}^m V(\rho(v_i))$, which is clearly finite-dimensional.
\end{proof}

\begin{lemma}\label{prop:reconstruction-cat-aux}
	With the notation above, $u: \coEnd(\omega) \to L$ in
	\eqref{eqn:u-comparison} is an isomorphism of coalgebras.
\end{lemma}
\begin{proof}
	We show that all right $L$-comodules $V$, together with all homomorphisms
	between them, lift naturally to right $\coEnd(\omega)$-comodules. For
	$\dim_{\Bbbk} V < \infty$, this is by definition---apply
	Definition--Proposition \ref{def:cogebra-L} to
	$\mathcal{A} = \catesubd{Comod}{\mathrm{f}}L$ and $\omega = U$; the
	general case then follows from Lemma \ref{prop:comodule-fd-union}.
	
	Now take $V=L$ with the right $L$-comodule structure determined by
	$\Delta$. The lifted $\coEnd(\omega)$-comodule structure corresponds to
	$\tilde{\Delta}: L \to L \dotimes{\Bbbk} \coEnd(\omega)$. This is the
	same as saying that the triangular part of the following diagram
	commutes:
	\[\begin{tikzcd}
		L \arrow[r, "{\tilde{\Delta}}"] \arrow[rd, "\Delta"'] & L \dotimes{\Bbbk} \coEnd(\omega) \arrow[r, "{\epsilon \otimes \identity}"] \arrow[d, "{\identity \otimes u}"] & \coEnd(\omega) \arrow[d, "u"] \\
		& L \dotimes{\Bbbk} L \arrow[r, "{\epsilon \otimes \identity}"'] & L,
	\end{tikzcd}\]
	and its square part clearly commutes as well. Denote the composite of the
	first row by $\gamma$; then $u\gamma = \identity_L$. It remains only to
	prove that $\gamma$ is surjective.
	
	Let $\rho: V \to V \dotimes{\Bbbk} L$ be a finite-dimensional right
	$L$-comodule. The commutative diagram in the definition
	\[\begin{tikzcd}
		V \arrow[r, "\rho"] \arrow[d, "\rho"'] & V \dotimes{\Bbbk} L \arrow[d, "{\identity \otimes \Delta}"] \\
		V \dotimes{\Bbbk} L \arrow[r, "{\rho \otimes \identity}"'] & V \dotimes{\Bbbk} L \dotimes{\Bbbk} L
	\end{tikzcd}\]
	is also equivalent to saying that $\rho$ is a homomorphism of
	$L$-comodules, provided that $V \dotimes{\Bbbk} L$ is equipped with the
	comodule structure coming from $L$ and $\Delta$. Now lift $V$ naturally
	to a $\coEnd(\omega)$-comodule, corresponding to
	$\tilde{\rho}: V \to V \dotimes{\Bbbk} \coEnd(\omega)$; if the comodule
	$V \dotimes{\Bbbk} L$ is lifted similarly, its structure corresponds to
	$\identity \otimes \tilde{\Delta}: V \dotimes{\Bbbk} L \to V \dotimes{\Bbbk} L \dotimes{\Bbbk} \coEnd(\omega)$. The comodule homomorphism
	$\rho$ also lifts to the $\coEnd(\omega)$ level, which is the same as
	saying that the square in the following diagram commutes:
	\[\begin{tikzcd}
		V \arrow[r, "\rho"] \arrow[d, "{\tilde{\rho}}"'] & V \dotimes{\Bbbk} L \arrow[d, "{\identity \otimes \tilde{\Delta}}"] & \\
		V \dotimes{\Bbbk} \coEnd(\omega) \arrow[r, "{\rho \otimes \identity}"' inner sep=0.8em] & V \dotimes{\Bbbk} L \dotimes{\Bbbk} \coEnd(\omega) \arrow[r, "{\identity \otimes \epsilon \otimes \identity}"' inner sep=0.8em] & V \dotimes{\Bbbk} \coEnd(\omega)
	\end{tikzcd}\]
	Since the composite of the second row is $\identity$, we obtain
	$\tilde{\rho} = (\identity_V \otimes \gamma)\rho$. One checks that the
	map $\alpha_V: V^\vee \dotimes{\Bbbk} V \to \coEnd(\omega)$ induced by
	$\tilde{\rho}$ therefore factors as
	\[ V^\vee \dotimes{\Bbbk} V \xrightarrow{\text{induced by $\rho$}} L \xrightarrow{\gamma} \coEnd(\omega). \]
	
	Recalling the construction of $\coEnd(\omega)$, as $V$ varies, the
	images $\Image(\alpha_V)$ generate all of $\coEnd(\omega)$. Thus
	$\gamma$ is surjective.
\end{proof}

As one might expect, if $L$ is a bialgebra, then \eqref{eqn:u-comparison} is
also an isomorphism of bialgebras. The required argument is entirely routine
and need not be spelled out.

\begin{definition}\label{def:Fib-cat}
	\index[sym1]{Fibk@$\cate{Fib}(\Bbbk)$, $\cate{Fib}^{\otimes}(\Bbbk)$}
	For the chosen field $\Bbbk$, define the category
	$\cate{Fib}(\Bbbk)$ as follows.
	\begin{itemize}
		\item Its objects are data $(\mathcal{A}, \omega)$, where
		$\mathcal{A}$ is a locally finite abelian category (and hence
		essentially small), and
		$\omega: \mathcal{A} \to \cate{Vect}_{\mathrm{f}}(\Bbbk)$ is a
		faithful exact functor.
		\item A morphism from $(\mathcal{A}, \omega)$ to
		$(\mathcal{A}', \omega')$ is a datum $(F, \alpha)$, where
		$F: \mathcal{A} \to \mathcal{A}'$ is a functor and
		$\alpha: \omega \rightiso \omega' F$ is an isomorphism of functors.
		Identities and composition are defined in the evident way.
	\end{itemize}
	Similarly, define the category $\cate{Fib}^{\otimes}(\Bbbk)$, with the
	following additional conditions:
	\begin{itemize}
		\item $\mathcal{A}$ in an object $(\mathcal{A}, \omega)$ must be
		equipped with a monoidal structure, and $\omega$ must be a monoidal
		functor;
		\item $F$ in a morphism must also be a monoidal functor, and $\alpha$
		must be an isomorphism of monoidal functors.
	\end{itemize}
\end{definition}

Observe that if $(\mathcal{A}, \omega)$ is an object of
$\cate{Fib}^{\otimes}(\Bbbk)$, faithfulness and exactness imply that
$\End_{\mathcal{A}}(\munit) \to \Bbbk$ is a monomorphism of
$\Bbbk$-algebras; thus in this case
$\End_{\mathcal{A}}(\munit) \rightiso \Bbbk$.

On the other hand, denote the category of $\Bbbk$-coalgebras (respectively,
$\Bbbk$-bialgebras) and their homomorphisms by $\Bbbk\dcate{coAlg}$
(respectively, $\Bbbk\dcate{biAlg}$).
\begin{itemize}
	\item There are evident functors
	$\cate{Fib}(\Bbbk) \to \Bbbk\dcate{coAlg}$ and
	$\cate{Fib}^{\otimes}(\Bbbk) \to \Bbbk\dcate{biAlg}$. Both send an
	object $(\mathcal{A}, \omega)$ to $\coEnd(\omega)$, while a morphism
	$(\mathcal{A}, \omega) \to (\mathcal{A}', \omega')$ induces a
	coalgebra or bialgebra homomorphism
	$\coEnd(\omega) \simeq \coEnd(\omega' F) \to \coEnd(\omega')$.
	\item There are also functors in the opposite direction: for a coalgebra
	or bialgebra $L$, take the corresponding datum
	$\left(\catesubd{Comod}{\mathrm{f}}L, U \right)$; on morphisms, for a
	homomorphism $\phi: L \to L'$, a right $L$-comodule $(M, \rho)$ becomes
	a right $L'$-comodule through
	$M \xrightarrow{(\identity \otimes \phi) \rho} M \otimes L'$.
	This gives
	$\left(\catesubd{Comod}{\mathrm{f}}L, U \right) \to \left(\catesubd{Comod}{\mathrm{f}}L', U' \right)$.
\end{itemize}

\begin{theorem}\label{prop:Tannaka-reconstruction}
	\index{reconstruction theorem}
	Let $\Bbbk$ be a field. The constructions above give two pairs of
	mutually quasi-inverse functors
	\[\begin{tikzcd}[row sep=tiny]
		\cate{Fib}(\Bbbk) \arrow[shift left, r] & \Bbbk\dcate{coAlg}, \arrow[shift left, l] \\
		\cate{Fib}^{\otimes}(\Bbbk) \arrow[shift left, r] & \Bbbk\dcate{biAlg}. \arrow[shift left, l]
	\end{tikzcd}\]
\end{theorem}
\begin{proof}
	Write the functors in the two directions as
	$Z: \star \leftrightarrows \star :Y$. For the first pair, the isomorphism
	$u$ in Lemma \ref{prop:reconstruction-cat-aux} gives
	$ZY \rightiso \identity$, while the equivalence $\overline{\omega}$ in
	Theorem \ref{prop:Tannaka-aux1} (i) gives
	$\identity \rightiso YZ$.
	
	For the second pair, since it is already known that when $L$ is a
	bialgebra, $u$ in \eqref{eqn:u-comparison} is also an isomorphism of
	bialgebras, we still have $ZY \rightiso \identity$ in the monoidal
	setting; meanwhile, $\identity \rightiso YZ$ is now the content of
	Theorem
	\ref{prop:Tannaka-aux1} (ii).
\end{proof}

We now add the duality introduced in \S\ref{sec:monoidal-dual} to the
framework of Reconstruction Theorem \ref{prop:Tannaka-aux1}. As usual, we
consider only the simple situation in which $B = \Bbbk$ is a field.

\begin{theorem}\label{prop:Tannaka-aux2}
	Under the assumptions of Theorem \ref{prop:Tannaka-reconstruction},
	suppose further that $(\mathcal{A}, \omega)$ is an object of
	$\cate{Fib}^{\otimes}(\Bbbk)$ and that $\mathcal{A}$ is both left and
	right rigid. Then $\coEnd(\omega)$ is a Hopf algebra.
\end{theorem}
\begin{proof}
	It suffices to show that the bialgebra $\coEnd(\omega)$ has an antipode.
	Use the left rigidity of $\mathcal{A}$ to define the functor
	$X \mapsto {}^* X$. For each $X \in \Obj(\mathcal{A})$, since monoidal
	functors preserve duals, we obtain
	$\omega({}^* X) \simeq \omega(X)^\vee$; since the monoidal structure on
	$\cate{Vect}_{\mathrm{f}}(\Bbbk)$ is symmetric, we also obtain
	$\omega(X) \simeq \omega({}^* X)^\vee$. Now consider
	\[ \phi \in \omega(X)^\vee \simeq \omega({}^* X), \quad x \in \omega(X) \simeq \omega({}^* X)^\vee , \]
	where $X \in \Obj(\mathcal{A})$ is arbitrary. In the notation of
	Convention \ref{con:coend-bracket}, both $[\phi \otimes x]$ and
	$[x \otimes \phi]$ are therefore elements of $\coEnd(\omega)$. We claim
	that
	\begin{equation}\label{eqn:Tannaka-aux2}
		\text{there is a linear map}\; S: \coEnd(\omega) \to \coEnd(\omega), \;\text{characterized by}\; S[\phi \otimes x] = [x \otimes \phi]. 
	\end{equation}

	For this, observe that for every morphism $f: X \to Y$ in
	$\mathcal{A}$, there are commutative diagrams
	\[\begin{tikzcd}
		\omega(X) \arrow[r, "{\omega(f)}"] \arrow[d, "\sim"' sloped] & \omega(Y) \arrow[d, "\sim" sloped] \\
		\omega({}^* X)^\vee \arrow[r, "{\omega({}^* f)^\vee}"' inner sep=0.6em] & \omega({}^* Y)^\vee
	\end{tikzcd} \quad \begin{tikzcd}
		\omega(Y)^\vee \arrow[r, "{\omega(f)^\vee}" inner sep=0.6em] \arrow[d, "\sim"' sloped] & \omega(X)^\vee \arrow[d, "\sim" sloped] \\
		\omega({}^* Y) \arrow[r, "{\omega({}^* f)}"'] & \omega({}^* X),
	\end{tikzcd}\]
	where all the vertical arrows come from the isomorphisms in the preceding
	paragraph. For $\psi \in \omega(Y)^\vee$ and $x \in \omega(X)$, there
	are corresponding elements $\psi' \in \omega({}^* Y)$ and
	$x' \in \omega({}^* X)^\vee$ such that, respectively,
	\[ \omega(f)^\vee (\psi) \leftrightarrow \omega({}^* f)(\psi'), \quad \omega(f)(x) \leftrightarrow \omega({}^* f)^\vee(x'). \]
	Substituting this into the construction and description of
	$\coEnd(\omega)$ in the proof of Proposition
	\ref{prop:cogebra-L-universal}, we see that $\delta_f$ there corresponds
	to $\delta_{{}^* f}$. Thus
	% Verified correction of the definite typo on source line 1578: the
	% source class is [\phi \otimes x], as in the definition of S above.
	$[\phi \otimes x] \mapsto [x \otimes \phi]$ indeed defines a unique
	linear map $S$, proving \eqref{eqn:Tannaka-aux2}.
	
	Our goal is to prove that $S$ is an antipode. The construction in the
	preceding paragraph can also be carried out with $X^*$ in place of
	${}^* X$, yielding a linear map
	$T: \coEnd(\omega) \to \coEnd(\omega)$; this map is likewise written in
	the form $[\phi \otimes x] \mapsto [x \otimes \phi]$. From
	${}^* (X^*) = X = ({}^* X)^*$, it is easy to see that
	$ST = \identity = TS$; in particular, $S$ is an isomorphism. This is the
	first requirement for an antipode.
	
	Next, denote the multiplication, comultiplication, unit, and counit
	homomorphisms of $\coEnd(\omega)$ by $\mu$, $\Delta$, $\eta$, and
	$\epsilon$, respectively. Choose an object $X$ and a basis
	$v_1, \ldots, v_n$ of the $\Bbbk$-vector space $\omega(X)$. Denote the
	dual basis of $\omega(X)^\vee$ by
	$\check{v}_1, \ldots, \check{v}_n$. Denote the morphisms for the left
	dual of $X$ by
	\[ \mathrm{ev}_X: {}^* X \otimes X \to \munit, \quad \mathrm{coev}_X: \munit \to X \otimes {}^* X. \]
	Their images under $\omega$ realize the duality between $\omega(X)$ and
	$\omega({}^* X)$. Recall that $(\cdot)^\vee$ is the dual in the symmetric
	monoidal category $\cate{Vect}_{\mathrm{f}}(\Bbbk)$, with no distinction
	between left and right; henceforth we may identify
	\begin{equation*}
		\omega(\munit) \simeq \omega(\munit)^\vee, \quad \omega(X \otimes {}^* X) \simeq \omega({}^* X \otimes X)^\vee,
	\end{equation*}
	and correspondingly identify $\omega(\mathrm{coev}_X)$ with
	$\omega(\mathrm{ev}_X)^\vee$.
	
	For every $\phi \in \omega(X)^\vee$ and $x \in \omega(X)$, the
	description of $\Delta$ (respectively, $\mu$) in
	\eqref{eqn:L-comult} (respectively, \eqref{eqn:L-tensor}) gives
	\begin{align*}
		\Delta([\phi \otimes x]) & = \sum_{i=1}^n [\phi \otimes v_i] \otimes [\check{v}_i \otimes x], \\
		\mu(S \otimes \identity)\Delta([\phi \otimes x]) & = \sum_{i=1}^n \mu\left( [v_i \otimes \phi] \otimes [\check{v}_i \otimes x] \right) \\
		& = \sum_{i=1}^n \left[ (v_i \otimes \check{v}_i) \otimes (\phi \otimes x) \right] \\
		& = \left[ \mathrm{coev}_{\omega(X)}(1) \otimes (\phi \otimes x) \right] \\
		& = \left[ \omega(\mathrm{coev}_X)(1) \otimes (\phi \otimes x) \right].
	\end{align*}
	However, \eqref{eqn:cogebra-L} implies that the following diagram
	commutes:
	\[\begin{tikzcd}
		\omega(\munit) \otimes \omega({}^* X \otimes X) \arrow[r, "{\omega(\mathrm{coev}_X) \otimes \identity}" inner sep=0.6em] \arrow[d, "{\identity \otimes \omega(\mathrm{ev}_X)}"'] & \omega(X \otimes {}^* X) \otimes \omega(X \otimes {}^* X), \arrow[d, "{[\cdot]}"] \\
		\omega(\munit) \otimes \omega(\munit) \arrow[d, "\sim"' sloped] \arrow[r, "{[\cdot]}"] & \coEnd(\omega) \\
		\Bbbk \arrow[ru, "\eta"'] & 
	\end{tikzcd}\]
	Therefore,
	\begin{gather*}
		\left[ \omega(\mathrm{coev}_X)(1) \otimes (\phi \otimes x) \right]
		= \eta\left( \omega(\mathrm{ev}_X)(\phi \otimes x) \right)
		\xlongequal{\text{\eqref{eqn:L-counit}}} \eta\epsilon\left( \phi \otimes x \right),
	\end{gather*}
	which proves $\mu(S \otimes \identity)\Delta = \eta\epsilon$.
	Similarly, $\mu(\identity \otimes S)\Delta = \eta\epsilon$. Thus $S$ is
	an antipode.
\end{proof}

\begin{corollary}\label{prop:Hopf-vs-dual}
	Under the bijection given by Theorem
	\ref{prop:Tannaka-reconstruction},
	\[\begin{tikzcd}
		\Obj(\cate{Fib}^{\otimes}(\Bbbk))/\simeq \arrow[shift left, r, "{1:1}"] & \Obj(\Bbbk\dcate{biAlg}) /\simeq , \arrow[shift left, l]
	\end{tikzcd}\]
	the category $\mathcal{A}$ in an object $(\mathcal{A}, \omega)$ on the
	left is both left and right rigid if and only if the object $L$ on the
	right is a Hopf algebra.
\end{corollary}
\begin{proof}
	The ``only if'' direction follows from Theorem \ref{prop:Tannaka-aux2},
	whereas the ``if'' direction was explained in Proposition
	\ref{prop:Hopf-module-dual}.
\end{proof}

In the proof of Theorem \ref{prop:Tannaka-aux2}, the left rigidity of
$\mathcal{A}$ plays the principal role; right rigidity is used only to
ensure that $S$ is an isomorphism. If the definition of a Hopf algebra
allows a noninvertible $S$, then the category of right comodules
$\catesubd{Comod}{\mathrm{f}}L$ is only left rigid; see the formulas in
Proposition \ref{prop:Hopf-module-dual}.
